\subsection{Серверная часть (Backend)}

Спроектированная модель данных служит основой для серверной части, обеспечивающей бизнес-логику и API для клиентского приложения. Серверная часть реализована на языке Rust с использованием фреймворка Axum. Архитектура следует паттерну разделения на слои: HTTP-обработчики (handlers), сервисы бизнес-логики (services), модели данных (models) и промежуточное программное обеспечение (middleware). На рисунке~\ref{backend-arch} представлена архитектура серверной части.

\begin{figure}
  \centering
  \includegraphics[width=0.95\textwidth,height=0.85\textheight,keepaspectratio]{backend-components.png}
  \caption{Архитектура серверной части}
  \label{backend-arch}
\end{figure}

\subsubsection{Маршрутизация и обработка запросов}

Все API-эндпоинты доступны по префиксу \texttt{/api/v1/}. Маршруты сгруппированы по доменным областям: животные, надои, воспроизводство, кормление, фитнес, отчёты, аналитика, настройки и другие. Общее количество обработчиков --- более 25 модулей. Примеры запросов и ответов приведены в приложении~\ref{app-api}.

Документация API генерируется автоматически с помощью библиотеки utoipa~\cite{utoipa} на основе аннотаций в коде и доступна через Swagger UI по адресу \texttt{/api/v1/docs}. Аннотации включают описание эндпоинтов, типы запросов и ответов, коды ошибок.

\subsubsection{Аутентификация и авторизация}

Аутентификация основана на \acrshort{jwt}~\cite{JWT} (JSON Web Token) с алгоритмом HS256. При входе в систему сервер выдаёт два токена: короткоживущий access token (15 минут) и долгоживущий refresh token (7 дней) для обновления access token без повторного ввода пароля.

Токены хранятся в HttpOnly, SameSite=Strict cookie, что предотвращает доступ к ним из JavaScript и обеспечивает защиту от CSRF-атак~\cite{OWASP-CSRF}. Отзыв токенов реализован через таблицу \texttt{revoked\_tokens}, хранящую JTI отозванных токенов. Просроченные записи удаляются ежечасно фоновым заданием.

Авторизация реализована через три уровня доступа, соответствующих типам Axum-экстракторов. Уровень \textbf{Claims} допускает любого аутентифицированного пользователя с валидным access token. Уровень \textbf{ClaimsAllowMustChange} допускает состояние обязательной смены пароля. Уровень \textbf{AdminGuard} ограничивает доступ пользователями с ролью \texttt{admin}.

\subsubsection{Промежуточное программное обеспечение (Middleware)}

Стек middleware обеспечивает безопасность, наблюдаемость и ограничение нагрузки. \textbf{CORS} настраивает кросс-доменный доступ~\cite{CORS} с конфигурируемыми источниками. \textbf{CSRF-защита} проверяет заголовок Origin для изменяющих состояние запросов. \textbf{Rate Limiting} ограничивает частоту запросов по IP-адресу (100 запросов за 60 секунд) с использованием Redis. \textbf{Gzip-сжатие} уменьшает объём передаваемых данных. \textbf{Логирование} обеспечивает трассировку HTTP-запросов через tracing с поддержкой структурированных логов в JSON-формате. \textbf{Метрики Prometheus} включают счётчики \texttt{http\_requests\_total} и гистограммы \texttt{http\_request\_duration\_seconds} с нормализацией маршрутов.

\subsubsection{Интеграция с Lely}

Модуль интеграции с Lely обеспечивает автоматическую двустороннюю синхронизацию данных с облачным API Lely Horizon. Синхронизация выполняется по настраиваемому таймеру (по умолчанию каждые 5 минут).

Архитектура модуля включает HTTP-клиент с аутентификацией по токену и автоматическим обновлением каждые 58 минут, а также повторными попытками с экспоненциальной задержкой. Планировщик синхронизации поддерживает параллельные задачи и инкрементальную синхронизацию по временным диапазонам (чанками по 7 дней). Маппинг данных использует кэш инвентарных номеров для разрешения Life Number Lely во внутренний идентификатор животного. Upsert-логика выполняет добавление или обновление записей через INSERT ON CONFLICT UPDATE. Учётные данные шифруются алгоритмом AES-256-GCM~\cite{AES-GCM}.

Синхронизируются 23 типа данных: животные, отёлы, осеменения, стельности, охоты, запуски, суточные надои, визиты на доение, качество молока, данные роботов, кормление, активности, жвачка, пастьба, быки-производители, переводы, кровные линии, типы и группы кормов, контакты, местоположения.
